% E4a/E4p plus the frozen validation supplementary archives; validation, not test.
\begin{table}[htbp]
\centering\small
\setlength{\tabcolsep}{3pt}
\bicaption[tab:component-results]{表}{FiLM-SIREN与FSC的五个开发阶段归档端点比较（44名被试，767方向）}{Tab.}{Validation comparison of FiLM-SIREN and FSC on five archived development endpoints (44 subjects, 767 directions)}
\begin{tabular}{lrr}
\toprule
指标$\downarrow$ & FiLM-SIREN E130 & FSC E190\\
\midrule
全域ERB & $0.8275 \pm 0.1602$ & $0.8044 \pm 0.1614$ \\
对侧$25^\circ$ ERB & $1.2308 \pm 0.1878$ & $1.2146 \pm 0.1887$ \\
频带ILD & $1.5739 \pm 0.1653$ & $1.4992 \pm 0.1523$ \\
ITD加权MAE（$\mu$s） & $15.5432 \pm 2.3330$ & $15.5224 \pm 2.3298$ \\
ITD最大绝对误差（$\mu$s） & $130.3267 \pm 39.6414$ & $130.2083 \pm 39.3381$ \\
\bottomrule
\end{tabular}
\par\smallskip
\parbox{\linewidth}{\footnotesize 数值为被试均值$\pm$样本标准差，两种模型均为三个成员的等权集成。两项ERB来自既有验证归档；频带ILD和两项ITD来自同一被试、方向域及冻结预测的补充验证归档。该表用于描述第二阶段完整训练链路的增量变化，不作为正文五项测试指标的替代。}
\end{table}
